\documentclass[../main]{subfiles}
\begin{document}
\ifSubfilesClassLoaded{\setcounter{page}{4}}{}
\section{Comunismo e famiglia}
\label{sec:09_comunismo_famiglia}

Nell’opera “L’origine della famiglia, della proprietà privata e dello Stato”, Friedrich Engels afferma: “Secondo la concezione materialistica, il fattore determinante nel corso della storia è, in ultima analisi, la produzione e la riproduzione della vita materiale. Questo aspetto, tuttavia, si manifesta in due forme distinte. Da un lato, la produzione dei mezzi di sussistenza – cibo, vestiario, alloggi e gli strumenti necessari a tali scopi; dall’altro, la produzione dello stesso uomo, la perpetuazione della specie. L’organizzazione sociale in cui vivono gli uomini di una determinata epoca storica e di un determinato paese è condizionata da entrambe queste forme di produzione: dal livello di sviluppo, da un lato, del lavoro e, dall’altro, della famiglia”.

Nelle sezioni precedenti, abbiamo sostenuto che, nel contesto del comunismo, la produzione dei beni essenziali è automatizzata e volta a soddisfare le esigenze dell’esistenza umana nella sua completezza. La società, agendo consapevolmente, plasma il proprio sviluppo, definisce le forme della sua organizzazione collettiva e, al contempo, espande i confini dell’attività libera. La famiglia, intesa come unità economica fondata sul carattere privato dell’appropriazione dei frutti del lavoro, vede declinare la sua rilevanza e, infine, scompare durante la fase di transizione. Come Engels osservò, «l’economia domestica privata si trasformerà in un settore del lavoro sociale». Abbiamo già evidenziato quanto questo processo consenta di risparmiare tempo prezioso per la società, alleviando gli individui dalle fatiche della vita quotidiana. Ciononostante, la socializzazione non si limita alla sfera del consumo.

La trasformazione più rilevante nell’ambito dei rapporti familiari risiede nell’evoluzione verso un modello di educazione collettiva dei figli. Tale processo, in atto da tempo, aveva già manifestato i suoi primi segnali anche nell’era capitalista. L’organizzazione domestica ha cominciato a “socializzarsi”, e da tempo numerosi aspetti della vita quotidiana sono diventati parte integrante del sistema economico pubblico. La preparazione dei pasti, ad esempio, si avvale prevalentemente di prodotti semilavorati; l’abbigliamento viene acquistato nei negozi e, sempre più spesso, la sua cura è affidata a servizi esterni, sebbene in passato tali attività fossero prerogative femminili, svolte all’interno dell’abitazione. Analogamente, l’uomo contemporaneo non è più gravato dalla necessità di trasportare l’acqua, spaccare la legna o costruire la dimora per la sua famiglia. Tutte queste incombenze sono state da tempo assorbite dal sistema produttivo pubblico e tecnologico.

L’educazione dei bambini non è ancora diventata un vero e proprio settore produttivo, nonostante l’esistenza e lo sviluppo attivo, per almeno due secoli, delle istituzioni pubbliche preposte. Da sempre, infatti, le famiglie più agiate hanno affidato i propri figli a tali istituzioni, ricorrendo anche a insegnanti privati, mentre i genitori si sono goduti la serenità della vita familiare. Anche nell’Unione Sovietica, uno dei primi traguardi di cui andavano fieri i lavoratori socialisti era l’organizzazione di asili nido e scuole materne presso le fabbriche. In altre parole, l’esigenza di affidare l’educazione dei bambini a persone competenti, al fine di guadagnare tempo libero da dedicare a ciò che si ritiene realmente utile e importante (e, per alcuni, questo non è forse proprio l’educazione dei figli?), è un desiderio condiviso sia dai benestanti che dagli indigenti, sebbene con alcune lievi differenze di indole, forse di carattere borghese o tendenzialmente poetico.Il capitalismo, fin dalla notte dei tempi, erode la struttura della famiglia tradizionale, e questo declino non affonda le sue radici in una crisi morale, bensì in dinamiche puramente economiche. Oggi, l'enfasi sulla promozione dei valori familiari risponde, da un lato, all'esigenza di controbilanciare i successi delle rivoluzioni socialiste, che hanno abolito lo sfruttamento del lavoro minorile, e, dall'altro, alla necessità, per il capitale, di creare un mercato per i beni e i servizi destinati all'infanzia e all'istruzione. Ma non sussiste una reale necessità storica e sociale di preservare la famiglia nella sua attuale configurazione di unità economica.

Di conseguenza, con l’evoluzione della società verso un sistema comunista, l’ambito dell’educazione infantile si trasformerà in un settore della produzione sociale. L’istruzione diventerà una vera e propria attività sociale, un’azione libera guidata dai principi della verità, del bene e della bellezza. Ciò comporta un aumento considerevole della responsabilità degli educatori, sia nei confronti dei bambini, sia nei confronti dell’intera società. Come affermava A.S. Makarenko: «Nel nostro sistema educativo, non devono verificarsi né catastrofi infantili, né insuccessi, né percentuali di errore, nemmeno espresse in centesimi!» In diverse sue opere, paragonava l’educazione ad altri ambiti, evidenziando come un fallimento in questo campo possa rivelarsi più pericoloso di qualsiasi altro. A proposito della famiglia, Anton Semënovič sottolineava che il compito dei genitori consiste nell’allevare individui degni di far parte della società, e che, per quanto riguarda la «qualità del contributo fornito», essi sono responsabili verso la comunità, ricevendo in cambio l’affetto e l’armonia dei legami familiari.Non tutti i genitori, tuttavia, possiedono le competenze necessarie per una corretta educazione dei figli. Inoltre, un numero crescente di bambini si trova privo di figure genitoriali, o vive in contesti familiari incompleti. Da qui l’esigenza di un impegno sociale più consapevole e strutturato nel campo dell’educazione, e della creazione di nuove forme di aggregazione familiare.

La «famiglia» comunista rappresenta una forma di aggregazione collettiva, all'interno della quale ogni individuo contribuisce alla creazione del tessuto sociale su una scala più ampia. Ne consegue che le dimensioni di tale «famiglia» possono variare, e non possono certo essere limitate alle attuali dimensioni. Queste dipendono da molteplici fattori storico-sociali, legati alle modalità di produzione, e di conseguenza saranno possibili diverse varianti, sia all'interno di una specifica epoca nello sviluppo della società comunista, sia nel corso della sua evoluzione storica. Poiché questa «famiglia», in quanto collettività primaria nella formazione dell'individuo, si fonda sulla natura sociale dell'attività umana e ne è direttamente determinata, ciò implica che anche il principio alla base della sua creazione sarà intrinsecamente sociale, e quindi qualitativamente diverso rispetto alle dinamiche attuali.Ciò che è in discussione è la dissoluzione del matrimonio moderno nella sua forma istituzionale, soprattutto perché le relazioni sessuali, l’amore e i legami matrimoniali e familiari rappresentano realtà profondamente distinte. Di conseguenza, una delle scoperte più significative e durature di A.S. Makarenko fu l’affermazione che la nuova “famiglia” deve integrare attivamente ogni individuo nel tessuto della vita sociale, coinvolgendolo come membro partecipe di un collettivo impegnato nella produzione, e non relegandolo al ruolo di mero consumatore di beni sociali o di semplice “oggetto” di un processo educativo.

Nel suo saggio “Sulla via di un Eros alato!”, Aleksandra Kollontaj analizzò i rapidi mutamenti – talvolta impercettibili a un occhio inesperto – che investirono questo settore nei primi anni della Grande Rivoluzione socialista d’Ottobre. Cercò, in particolare, di individuare e anticipare le forme di relazione personale, o meglio, d’amore, che avrebbero preso piede con l’affermazione definitiva della proprietà pubblica sulla privata, con lo sviluppo di nuove istituzioni sociali e con l’introduzione di strumenti volti a facilitare le incombenze legate alla maternità, trasformandola progressivamente in una componente essenziale della produzione sociale.

La nuova realtà, che si affermò (non solo nell'URSS), offrì alle donne opportunità di emancipazione impensabili nel resto del mondo, concedendo loro il tempo non soltanto per coltivare l'amore e il corteggiamento, ma anche per meditare sul concetto stesso di amore. Queste nuove relazioni, già in nuce di stampo comunista, sono descritte anche nella letteratura dell'epoca dell'industrializzazione, sia nella prosa di Aleksandra Kollontaj – ne «Vasilisa Malygina» e «Grande amore» – sia, per esempio, nel romanzo «Il secondo giorno» di Il'ja Erenburg.

L'affermazione della convivenza al di fuori del matrimonio legalmente sancito – la convivenza di fatto, divenuta prassi diffusa negli anni Venti e Trenta, mentre la registrazione dell'unione assumeva il significato opposto, ovvero l'indice di una mancanza di fiducia reciproca – in un'epoca in cui prevaleva ancora il modello della famiglia patriarcale, rappresentò senza dubbio un'anticipazione dei tempi. Più che una semplice rivoluzione, fu una trasformazione radicale nella concezione che le persone avevano di sé, dell'amore, dell'unione, della famiglia e persino del divorzio. Non più il matrimonio inteso come un'entità esclusiva per due individui, a beneficio di entrambi, ma piuttosto una realtà familiare direttamente influenzata dalle esigenze di una società in rapida espansione e trasformazione, e dalle dinamiche del contesto lavorativo più immediato: in questo risiede la vera essenza dell'amore, che manifesta una tendenza intrinsecamente comunitaria.

In quanto nucleo affettivo e, al contempo, teatro delle responsabilità che definiscono l’individuo nella società capitalista, la famiglia rappresenta un “territorio ad alto rischio”. Tra le sue mura, i coniugi sono chiamati a confrontarsi con le dinamiche predefinite dai rapporti di proprietà privata, con le questioni legate alla ricchezza, al guadagno, alle spese e alla banale quotidianità. Ma il comunismo, nella sua essenza, mira proprio all’eliminazione di tali problematiche.

La dinamica del declino, sia per quanto concerne la famiglia tradizionale sia per l’istituzione dello Stato stesso, è un fenomeno complesso. Tuttavia, le soluzioni a metà strada proposte dal socialismo si sono rivelate inefficaci: dagli alloggi comunitari, infatti, studenti e giovani professionisti finivano comunque per ritrovarsi in appartamenti, all’interno di nuclei familiari, pur continuando a fare riferimento alle vecchie strutture, sebbene sotto una nuova forma di assistenza sociale. Le motivazioni che avrebbero dovuto portare al superamento delle forme di organizzazione sociale preesistenti si sono perdute in una società in cui le relazioni monetarie e commerciali hanno riacquistato una posizione di predominio. Dopo questo tentativo di promuovere la comunità, è seguito un inevitabile ritorno alla struttura patriarcale, anch’esso, apparentemente, di natura transitoria. Nuovamente, l’attenzione si è concentrata prevalentemente sui genitori, piuttosto che sull’intero gruppo sociale o sulla comunità, nell’ambito dell’educazione dei figli. In realtà, si stavano delineando non solo i “figli della disciplina” (la guerra aveva apportato delle modifiche), ma anche i figli delle comunità.

In questi «figli della comunità» si manifestava uno degli elementi cruciali per l’affermazione del comunismo. Nel processo di transizione verso il comunismo, i bambini devono superare la distinzione tra «noi» e «loro». In un vasto collettivo, dove vengono risolti i problemi relativi alla «quantificazione» del contributo lavorativo individuale, non devono esistere bambini «altrui» (cioè bambini di cui qualcun altro, e non io, dovrebbe prendersi cura).

In un contesto sociale ampio, dove uomini e donne, adulti e bambini, interagiscono, la distinzione tra figura paterna e materna si attenua, e con essa svanisce l’angoscia di una presunta «fine del mondo». L’indebolimento dei legami genitoriali, sostituito da una rete di relazioni più estesa, quella dei «vicini», permette ai bambini di trascendere il gioco limitato alla «sabbia della proprietà privata». L’abbandono del modello tradizionale di coppia apre la strada a un futuro liberato da schemi obsoleti, residui del passato che persistono nell’inconscio e nella coscienza, non solo dei bambini, ma anche degli adulti. La trasformazione delle dinamiche di coppia, all’interno di un gruppo coeso e solidale, cessa di rappresentare un evento traumatico, soprattutto quando l’amore non è più il collante tra i membri della coppia.In linea generale, e in assenza di ostacoli di natura economica o sociale, è solo ciò che riguarda la sfera più intima e personale che può acquisire autenticità, pur senza opporsi alla società, ma anzi liberandosi dalle costrizioni economiche, legali e morali.

La nuova concezione di «famiglia» deve evolvere in una forma di impegno comunista, in cui ogni suo membro, compresi i bambini, sia attivamente coinvolto e connesso alla società, agendo come soggetto partecipe.
\end{document}
